\begin{proof}[Proof of \cref{obj:lem:calibration-consequences}]
Because \(n\ge N_0\), the calibration cutoff gives
\[
\alpha_0L\ge2,\qquad
4M(n)^2\le m_1,\qquad
\frac{4M(n)}{m_1}\le \frac{3B(n)}4.
\]
In particular \(L>0\). The identity \(L=\log(en)=1+\log n=1+\ell\) is valid for this positive \(n\). Since \(L\ge240\), we have
\[
\ell=L-1>0,\qquad \ell\le L,\qquad L=1+\ell\le2\ell. % lean: large_calibration_bounds
\]

The definition \(M(n)=\max\{2,\lfloor\alpha_0L\rfloor\}\) gives \(M\ge2\). Hence \(4M^2\le m_1\) implies \(0<m_1\). For the deterministic second split,
\[
m_1=n-\lfloor n/2\rfloor,
\]
and therefore \(n\le2m_1\). % lean: splitSize_one_eq

Since \(\alpha_0L\ge2\), the maximum in the definition of \(M(n)\) is attained by the floor term:
\[
M=\lfloor\alpha_0L\rfloor.
\]
Using \(\alpha_0\le1\) and \(L>0\),
\[
M\le \alpha_0L\le L. % lean: large_calibration_bounds
\]

By the definition \(B(n)=4096L/m_1\) and the bound \(n\le2m_1\),
\[
B=\frac{4096L}{m_1}
\le
\frac{8192L}{n}. % lean: large_calibration_bounds
\]

Finally we verify the calibrated growth inequality. Since \(\alpha_0\le(32\log6)^{-1}\) and \(M\le\alpha_0L\),
\[
4M\log6\le \frac{L}{8},
\qquad\text{so}\qquad
6^{4M}\le \exp(L/8).
\]
Also \(L\ge240\) gives
\[
\log L=\log 12+\log(L/12)
\le 11+(L/12-1)\le L/8,
\]
and hence
\[
L^6=\exp(6\log L)\le \exp(3L/4).
\]
Multiplying,
\[
6^{4M}L^6\le \exp(7L/8)\le \exp(L-1)=n.
\]
Since \((6^{2M})^2=6^{4M}\), this is exactly
\[
(6^{2M})^2L^6\le n. % lean: calibrated_polynomial_log_growth
\]
The displayed conclusions follow. % lean: large_calibration_bounds
\end{proof}
